\section{Fourier Sine Series}

\subsubsection*{Motivation and Intuition}

If a function is odd, then it is symmetric through the origin. Since sine functions are odd, an odd function can be represented using only sine terms.

\begin{conceptbox}
    If $f(x)$ is odd on $[-L,L]$, then
    \[
        a_0=0,\qquad a_n=0
    \]
    for all $n$, and the Fourier series becomes
    \[
        f(x)=\sum_{n=1}^{\infty}
        b_n\sin\left(\frac{n\pi x}{L}\right),
    \]
    where
    \[
        b_n=\frac{2}{L}\int_0^L
        f(x)\sin\left(\frac{n\pi x}{L}\right)\,dx.
    \]
\end{conceptbox}

\subsubsection*{Geometric and Physical Interpretation}

A Fourier sine series represents an odd periodic extension of a function defined on $0<x<L$.

\[
    f(-x)=-f(x)
\]

This is useful when boundary values are naturally zero, such as vibrating strings fixed at both ends.

\subsubsection*{Worked Example}

\begin{examplebox}
    Find the Fourier sine series of
    \[
        f(x)=1,\qquad 0<x<\pi.
    \]

    For a half-range sine expansion on $(0,\pi)$, we use
    \[
        f(x)=\sum_{n=1}^{\infty}b_n\sin(nx),
    \]
    where
    \[
        b_n=\frac{2}{\pi}\int_0^\pi 1\cdot \sin(nx)\,dx.
    \]

    Thus,
    \[
        b_n=\frac{2}{\pi}\int_0^\pi \sin(nx)\,dx.
    \]

    Compute the integral:
    \[
        \int \sin(nx)\,dx=-\frac{\cos(nx)}{n}.
    \]

    Therefore,
    \[
        b_n=
        \frac{2}{\pi}
        \left[
            -\frac{\cos(nx)}{n}
            \right]_0^\pi.
    \]

    Evaluate:
    \[
        b_n=
        \frac{2}{\pi}
        \left[
            -\frac{\cos(n\pi)}{n}
            +
            \frac{\cos(0)}{n}
            \right].
    \]

    Since
    \[
        \cos(n\pi)=(-1)^n,\qquad \cos(0)=1,
    \]
    we get
    \[
        b_n=
        \frac{2}{n\pi}
        \left[1-(-1)^n\right].
    \]

    If $n$ is even, then $(-1)^n=1$, so
    \[
        b_n=0.
    \]

    If $n$ is odd, then $(-1)^n=-1$, so
    \[
        b_n=\frac{4}{n\pi}.
    \]

    Thus,
    \[
        1=
        \frac{4}{\pi}
        \left[
            \sin x+\frac{1}{3}\sin(3x)+\frac{1}{5}\sin(5x)+\cdots
            \right],
        \qquad 0<x<\pi.
    \]
\end{examplebox}

\subsubsection*{Engineering Insight}

\begin{noteBox}
    Fourier sine series naturally appear in vibration and wave problems where the endpoints are fixed. For example, a string fixed at both ends has displacement zero at the boundaries, and sine functions automatically satisfy this condition.
\end{noteBox}

\subsubsection*{Common Mistakes and Notes}

\begin{conceptbox}
    A Fourier sine series represents an odd extension. At $x=0$ and $x=L$, the sine series evaluates to zero. Therefore, if the original function has nonzero endpoint values, the sine series may still represent the function on the open interval $(0,L)$ but not necessarily at the endpoints.
\end{conceptbox}